Найти поля частных для колец:

а) $\Z$; \;\;
б) $\Q$; \;\;
в) $\Z \left[\frac{1}{2}\right]$\;\;
г) $\Z[x]$;\;\;
д) $\Z[i]$;\;\;
е) $\R$;\;\;
ж) $\Z[\omega]$;\;\;
з) $\Q[x]$.
\bigbreak
Будем пользоваться свойствами из билета 27 на хор:

1) Если $F$ -- поле, то $Quot(F) \cong F$.

2) Если $K \subset L \subset Quot(K)$, то $Quot(K) \cong Quot(L)$.
\bigbreak
а) $Quot(\Z) \cong \Q$ \;\; (по определению поля частных)

б) $Quot(\Q) \cong \Q$ \quad(так как $\Q$ -- поле)

в) $Quot(\Z \left[ \frac{1}{2} \right]) \cong \Q$ \quad
(так как $\Z \subset \Z \left[ \frac{1}{2} \right] \subset \Q$)

г) $Quot(\Z[x]) = \left\{ \frac{f(x)}{g(x)} \;\Big| \; 
g(x) \neq 0, f(x), g(x) \in \Z[x] \right\} \cong \Z(x)$ \quad
(это просто определение $\Z(x)$)

д) $Quot(\Z[i]) \cong \Q[i]$.\;\; У нас получится отношение вида
$\frac{a + bi}{c + di}$, которое можно домножить на сопряженное и получить
элемент из $\Q[i]$.

е) $Quot(\R) \cong \R$ \quad(так как $\R$ -- поле)

ж) $Quot(\Z[w]) \cong \Q[w]$ \quad(по аналогии с пунктом д)

з) $Quot(\Q[x]) \cong \Q(x) \cong \Z(x)$